\chapter{Data model}
\label{ch:datamodel}

\section{Objects, ensembles and the domain}

An application creates one \emph{ensemble} per decomposition. The ensemble owns the domain ---
the global box, its periodicity and the assignment of sub-boxes to ranks --- and the list of
objects:
\begin{lstlisting}[numbers=none]
MFTL_Ensemble *MFTLEnsembleCreate(MPI_Comm comm, int cuts[3], double global_min[3],
                                  double global_max[3], bool is_periodic[3],
                                  int initial_object_capacity);
MFTL_Object *MFTLObjectCreate(int global_id);
int  MFTLEnsembleAddObject(MFTL_Ensemble *ensemble, MFTL_Object *object);
void MFTLEnsembleDestroy(MFTL_Ensemble **p_ens);
\end{lstlisting}
\fn{MFTLEnsembleCreate} is collective. It duplicates the communicator, checks that
$\texttt{cuts[0]}\cdot\texttt{cuts[1]}\cdot\texttt{cuts[2]}$ equals the number of ranks, splits the
box uniformly into that many sub-boxes, builds a Cartesian communicator and a
distributed-graph communicator over the (up to 26) neighbouring sub-boxes, and registers the
MPI datatypes of the library. \fn{MFTLEnsembleCreateRectilinear} accepts explicit slab
coordinates per axis and an explicit slab-to-rank map; \multi{} uses it to align the ensemble's
decomposition with the grid of its front-tracking overlap mesh.
\fn{MFTLEnsembleSetDecomposition} replaces the decomposition of an existing ensemble (the objects
must be redistributed afterwards).

An object is created empty with its object id and filled by a generator, a reader or a
builder. \fn{MFTLEnsembleAddObject} links it to the ensemble's domain. The object list must be
the same on every rank: every rank creates and adds the same objects in the same order, even
a rank that stores no part of them. \fn{MFTLEnsembleDestroy} destroys the objects with the
ensemble; \fn{MFTLEnsembleRemoveObject} detaches an object without destroying it, and
\fn{MFTLEnsembleTransferObject} moves an object between two ensembles with the same
decomposition.

\section{Vertices, triangles and half-edges}

Each object stores this rank's part of its mesh in three flat arrays: \texttt{vertices[]},
\texttt{triangles[]} and \texttt{half\_edges[]} with counts \texttt{n\_vertices},
\texttt{n\_triangles} and \texttt{n\_half\_edges}. References between entities are array indices
on the rank (\emph{local indices}). Every entity also has a \texttt{global\_id} that is unique
within its object and identical on every rank that stores a copy; the maps
\texttt{vertex\_lookup}, \texttt{triangle\_lookup} and \texttt{halfedge\_lookup} (\texttt{khash}
hash maps) translate a global id into the local index. Local indices change whenever the
arrays are rebuilt (redistribution, compaction after remeshing); global ids do not.

\begin{figure}[ht]
\centering
\begin{tikzpicture}[scale=2.2,
  he/.style={-{Stealth[length=2.2mm]},thick},
  tw/.style={-{Stealth[length=2.2mm]},thick,blue!70!black},
  v/.style={circle,fill=black,inner sep=1.3pt}]
  \coordinate (A) at (0,0); \coordinate (B) at (2,0); \coordinate (C) at (1,1.35); \coordinate (D) at (1,-1.35);
  \draw[fill=gray!8] (A)--(B)--(C)--cycle; \draw[fill=blue!5] (A)--(D)--(B)--cycle;
  \node[v,label=left:$a$] at (A) {}; \node[v,label=right:$b$] at (B) {};
  \node[v,label=above:$c$] at (C) {}; \node[v,label=below:$d$] at (D) {};
  \coordinate (G) at (1,0.45); % centroid of T
  \draw[he] ($(A)!0.15!(B)!0.18!(G)$) -- node[above]{\small $h$} ($(A)!0.85!(B)!0.18!(G)$);
  \draw[he] ($(B)!0.15!(C)!0.18!(G)$) -- ($(B)!0.85!(C)!0.18!(G)$);
  \draw[he] ($(C)!0.15!(A)!0.18!(G)$) -- ($(C)!0.85!(A)!0.18!(G)$);
  \node[font=\small] at (1.95,0.85) {\texttt{next}$(h)$};
  \node[font=\small] at (0.05,0.85) {\texttt{prev}$(h)$};
  \coordinate (H) at (1,-0.45); % centroid of T'
  \draw[tw] ($(B)!0.15!(A)!0.18!(H)$) -- node[below]{\small \texttt{twin}$(h)$} ($(B)!0.85!(A)!0.18!(H)$);
  \node at (1,0.62) {$T$}; \node at (1,-0.75) {$T'$};
\end{tikzpicture}
\caption{Half-edges of two neighbouring triangles. The half-edge $h$ runs from its origin $a$
to its target $b$ (\texttt{vertex\_idx} is the index of $b$); \texttt{next\_idx} cycles through
the three half-edges of $T$ in winding order; the twin runs $b\to a$ in the neighbouring
triangle $T'$. The origin of $h$ is the target of \texttt{prev}$(h)$, so it is found without
the twin.}
\label{fig:halfedge}
\end{figure}

\paragraph{Vertex} (\texttt{MFTL\_Vertex}):
\begin{itemize}
  \item \texttt{pos[3]}: the position, \emph{wrapped} into $[\min,\max)$ on periodic axes;
    \texttt{img[3]}: the periodic image counters (Chapter~\ref{ch:periodicity});
  \item derived geometry: \texttt{normal}, \texttt{meancurvature} ($2H$),
    \texttt{gaussiancurvature}, \texttt{area} (mixed Voronoi area), \texttt{roughness\_mean},
    \texttt{roughness\_variance}, and \texttt{curv\_extrap} (extrapolated curvature of a
    boundary vertex);
  \item \texttt{owner\_proc}, \texttt{is\_ghost}, \texttt{global\_id};
  \item \texttt{out\_he\_idx}: the local index of one outgoing half-edge --- the entry point of
    the one-ring walk. On a boundary vertex it is the outgoing boundary half-edge;
  \item \texttt{flags}: \texttt{MFTL\_FLAG\_*} and \texttt{MFTL\_VFLAG\_*} bits (remeshing
    marks, sharp features, curvature validity, boundary corners).
\end{itemize}

\paragraph{Triangle} (\texttt{MFTL\_Triangle}):
\begin{itemize}
  \item \texttt{he\_idx}: the local index of one of its three half-edges;
  \item cached \texttt{normal}, \texttt{area} and \texttt{center} (set by
    \fn{MFTLComputeTriangleProperties}, see \S\ref{sec:triprops});
  \item \texttt{owner\_proc}, \texttt{is\_ghost}, \texttt{global\_id}, \texttt{flags};
  \item \texttt{data}: an opaque payload of \texttt{data\_size} bytes, and
    \texttt{surface\_scalar}, a conserved concentration per unit area (Chapter~\ref{ch:data}).
\end{itemize}

\paragraph{Half-edge} (\texttt{MFTL\_HalfEdge}): a directed edge of exactly one triangle.
\begin{itemize}
  \item \texttt{vertex\_idx}: the vertex at its \emph{end} (target);
  \item \texttt{next\_idx}: the next half-edge of the same triangle in winding order; the three
    half-edges of a triangle form a cycle;
  \item \texttt{triangle\_idx}: its triangle;
  \item \texttt{twin\_idx} and \texttt{twin\_gid}: the opposite half-edge. Together they encode
    the \emph{twin state} LOCAL, REMOTE or BOUNDARY (Chapter~\ref{ch:topology}). They are
    never read directly: the predicates \fn{MFTLHeIsLocal}, \fn{MFTLHeIsRemote},
    \fn{MFTLHeIsBoundary}, \fn{MFTLHeBoundaryTag}, \fn{MFTLHeTwinGid} and \fn{MFTLHeEdgeId} are
    the only sanctioned readers (the script \texttt{tools/check\_twin\_gid\_readers.sh} enforces
    this in the library sources);
  \item \texttt{global\_id}, \texttt{flags}, and \texttt{boundary}, a cache of
    \fn{MFTLHeIsBoundary}.
\end{itemize}

\section{Navigation}

\texttt{mftl\_topology.h} provides inline accessors. With a half-edge $h$:
\fn{MFTLHeTarget} and \fn{MFTLHeOrigin} give its end and start vertex, \fn{MFTLHeNext} and
\fn{MFTLHePrev} the other half-edges of its triangle, \fn{MFTLHeTriangle} its triangle and
\fn{MFTLHeTwin} its stored twin. \fn{MFTLHeTwin} returns \texttt{NULL} when the twin is not
stored on this rank, which happens both at a physical boundary and at an MPI rim;
\texttt{NULL} must therefore never be taken to mean ``boundary''. For a triangle,
\fn{MFTLGetTriangleHalfedges} and \fn{MFTLGetTriangleVertices} return its half-edges and
vertices in winding order, and \fn{MFTLGetTriangleNeighbors} the triangles across its edges.

\paragraph{One-ring traversal.} The triangles around a vertex $v$ are visited through its
outgoing half-edges: from an outgoing half-edge $h$, the next one is
$\texttt{next}(\texttt{twin}(h))$. Seen from the side into which the triangle normals point
(from outside an outward-wound body, where the triangles appear counter-clockwise), this step
turns \emph{clockwise} around $v$, and $\texttt{twin}(\texttt{prev}(h))$ turns
counter-clockwise (Figure~\ref{fig:onering}). The walk stops where a half-edge has no stored
twin. \fn{MFTLGetVertexOneRing}, \fn{MFTLGetVertexOneRingVertices} and
\fn{MFTLGetVertexOneRingHalfedges} first rewind counter-clockwise to such a stop (unless the caller asserts a closed fan with \texttt{is\_inner}), so an
open fan is returned from one end to the other. \fn{MFTLVertexRing} classifies the fan of a
vertex as
\begin{itemize}
  \item \texttt{MFTL\_RING\_CLOSED}: an interior vertex whose whole fan is stored;
  \item \texttt{MFTL\_RING\_OPEN\_COMPLETE}: a boundary vertex; the fan runs from an incoming to
    an outgoing boundary half-edge;
  \item \texttt{MFTL\_RING\_INCOMPLETE}: the walk reaches a REMOTE half-edge (the fan is cut by
    an MPI rim) or the fan is malformed.
\end{itemize}

\begin{figure}[ht]
\centering
\begin{tikzpicture}[scale=1.6,
  he/.style={-{Stealth[length=2.2mm]},thick},
  v/.style={circle,fill=black,inner sep=1.2pt}]
  \coordinate (V) at (0,0);
  \foreach \i in {0,...,5} { \coordinate (P\i) at ({60*\i}:1.6); }
  \foreach \i/\j in {0/1,1/2,2/3,3/4,4/5,5/0} { \draw[fill=gray!6] (V)--(P\i)--(P\j)--cycle; }
  \foreach \i in {0,...,5} { \node[v] at (P\i) {}; }
  \node[v] at (V) {}; \node[below left] at (V) {$v$};
  \node[right] at (P0) {$P_0$}; \node[below left] at (P5) {$P_5$}; \node[above right] at (P1) {$P_1$};
  % seen from the normal side: triangles are counter-clockwise, so h = v->P0 lies in (v,P0,P1)
  \draw[he,red!70!black] ($(V)!0.12!(P0)+(90:0.08)$) -- node[above,pos=0.6]{\small $h$} ($(V)!0.85!(P0)+(90:0.08)$);
  % twin(h) = P0->v lies in (v,P5,P0)
  \draw[he,blue!70!black] ($(P0)!0.15!(V)+(-90:0.08)$) -- node[below,pos=0.35]{\small \texttt{twin}$(h)$} ($(P0)!0.88!(V)+(-90:0.08)$);
  % next(twin(h)) = v->P5, also in (v,P5,P0)
  \draw[he,red!70!black] ($(V)!0.12!(P5)+(30:0.08)$) -- ($(V)!0.85!(P5)+(30:0.08)$);
  \node[red!70!black,right] at ($(P5)+(0.05,0.1)$) {\small \texttt{next}(\texttt{twin}$(h)$)};
  \draw[-{Stealth},dashed] (-8:0.65) arc (-8:-52:0.65);
\end{tikzpicture}
\caption{The one-ring of $v$, seen from the side the normals point to (triangles
counter-clockwise). Red: the outgoing half-edges $h=v\to P_0$ and $\texttt{next}(\texttt{twin}(h))
=v\to P_5$; blue: $\texttt{twin}(h)$. The walk $h\mapsto\texttt{next}(\texttt{twin}(h))$ turns
clockwise around $v$.}
\label{fig:onering}
\end{figure}

\section{Flattened connectivity and neighbourhoods}

Some kernels (the LBVH, the neighbourhood search, the remesher's fits) use compact
connectivity arrays built by \fn{MFTLFlattenConnectivity}:
\texttt{flat\_tri\_verts[3t+j]} is the target vertex of the $j$-th half-edge of triangle $t$,
\texttt{flat\_tri\_neighbors[3t+j]} the triangle across that half-edge ($-1$ if its twin is
not stored), and \texttt{vert\_offsets}/\texttt{vert\_tris} a CSR map from a vertex to its
incident stored triangles. The generators, readers, builders and the redistribution rebuild
them; an application that calls the topological operators of the remesher directly calls
\fn{MFTLFlattenConnectivity} itself afterwards. \fn{MFTLGetNeighborhood} gathers vertices and triangles
breadth-first on this incidence graph (one step from a vertex gives its triangles, two steps
also their vertices).

\section{Owned entities first}

After a redistribution on more than one rank the owned entities of every array come first
and the ghosts follow; \texttt{n\_local\_vertices}, \texttt{n\_local\_triangles} and
\texttt{n\_local\_half\_edges} are then the lengths of the owned prefixes. These counts are set
only by a multi-rank redistribution, the builders (soups, \fn{MFTLBuildTopologyFromTriples}), the
staged copy \fn{MFTLObjectAppendTriangleCopy} and \fn{MFTLObjectClear};
generators, remeshing and a single-rank redistribution (which does nothing) do not update
them. Code that must work in every state therefore loops over the stored arrays and skips
entities with \texttt{is\_ghost} set --- the library itself identifies owned entities this way,
and so do all the examples in this manual.

\section{Global counts and id counters}

\texttt{global\_n\_triangles} and \texttt{global\_n\_vertices} of an object are the numbers of owned
entities summed over all ranks; they are refreshed by the collective
\fn{MFTLEnsembleUpdateGlobalCounts} (all objects) or \fn{MFTLObjectUpdateGlobalStats} (one
object). The counters \texttt{next\_vertex\_gid}, \texttt{next\_triangle\_gid} and
\texttt{next\_halfedge\_gid} hold the next unused global ids and are identical on every rank;
the parallel remesher hands out disjoint per-rank blocks of ids and re-synchronises the
counters afterwards (\fn{MFTLSyncEnsembleGidCounters}).
